
\subsection{Definition and scope}
\textbf{Trustworthy Premarital Health Education (\tphe)} is a proposed governance architecture for
premarital education that aims to increase each participant's control over marriage-related health
decisions through five rules, a stop--assess--refer--record pathway, repairable documentation, and
prospectively testable outcomes. \emph{Trustworthy} names a condition the programme must earn through
observable practice; \emph{premarital} names the transition point, not one legal or religious form;
\emph{health education} spans physical, psychological, sexual and reproductive, relational,
legal-protection, and help-seeking literacy without turning educators into clinicians, lawyers,
investigators, or religious adjudicators; \emph{infrastructure} means repeatable roles, boundaries,
records, referral relationships, training, and accountability. \tphe{} is placed upstream, but when it
encounters existing violence or coercion its function changes to identification and safe routing; it
connects to competent secondary and tertiary responses without providing them.

\tphe{} is \emph{not}: a diagnostic instrument; a substitute for clinical, social, legal, emergency, or
survivor-advocacy services; couple therapy or reconciliation counselling; a promise to prevent
violence; a requirement that any participant leave or proceed; a universal Islamic ruling or a
culturally neutral model; evidence that any organisation's programme is effective; an autonomous AI
counselling system; or a checklist detached from training, referral capacity, and governance.

\textbf{From six safeguards to five rules.} The accepted abstract names six safeguards that protect
decision control: epistemic humility, cultural intelligibility, affected-person visibility, dignity
and consent, bounded authority and referral, and repairability. Below, these are operationalised as
five rules, R1--R5 (Table~\ref{tab:rules}), because epistemic humility and cultural intelligibility
answer the same audit question --- whose authority is this claim carrying, and where does that
authority end --- and are therefore implemented together as R1 (bounded authority of knowledge and
culture). Affected-person visibility maps to R2; dignity and consent to R3; bounded authority and
referral splits across R1 and R4; repairability maps to R5. Every safeguard named in the accepted
abstract survives this operationalisation; the change from six to five is a regrouping of duties into
auditable rules, not a retraction of any safeguard, and the same mapping applies wherever the two
vocabularies (``six safeguards'' and ``five rules'') appear elsewhere in this paper.

\subsection{The model as a continuous line}
The programme theory does not treat the classroom day as the intervention. It treats the whole
trajectory from first contact through the day of check-in to whatever follows as one governed line,
on which the mechanism, five rules, and one test apply throughout. Three layers mark points on that
line; only the middle layer, the day itself, is part of the accepted programme.

\subsubsection*{Layer 0 --- signals from first contact (proposed extension)}
Before enrolment, course-facing material can state plainly what fear, control, pressure, or unsafe
participation look like, that pausing or refusing is a normal outcome, and where private help is
available. An anonymous scored self-check --- a form, activity, or game, usable at any point on the
line --- lets a person see their own result and routes without submitting a name or a narrative unless
they choose to. No identity is required to complete it, and nothing is sent onward unless the person
acts on it. The nearest evidence family is decision-aid research, not a self-check validated for this
population: a Cochrane review of 209 decision-aid trials found consistent gains in knowledge and
values-congruent choice \cite{ref14}, while the closest content-matched online self-efficacy tool for
women experiencing abuse, I-DECIDE, found \emph{no} improvement in self-efficacy at 6 or 12 months in a
randomised controlled trial \cite{ref119,ref120}. That null result is reported here, not
set aside, because it is the closest outcome-domain precedent this proposed layer has. Because
access to a shared or partner-monitored device is itself a disclosure risk independent of what is
entered, any pilot implementation must leave no discoverable trace on a shared device (no login, no
tool-specific browser or app history, a quick-exit function) and must state this limitation to
users rather than implying anonymity of content is sufficient.

\subsubsection*{Layer 1 --- the day as check-in}
The day of delivery is reframed as a check-in at which a participant verifies in person whether the
declared protections are in place: a private slot the partner cannot see or control, a professional
interpreter rather than a family member, and an assessor with no stake in whether the wedding
proceeds. Attendance on this day evidences that the space exists; it evidences nothing about whether
the programme changed anyone's decisions, knowledge, or safety. Five rules govern conduct on the whole
line and are auditable at any point, not only on the day. Table~\ref{tab:rules} states each rule's
duty, one observable trace, and one falsifier.

\begin{table*}[t]
\centering\footnotesize
\caption{The five rules (R1--R5), one observable trace and one falsifier each. Programme theory;
none of these traces has yet been measured against an outcome.}
\label{tab:rules}
\begin{tabularx}{\textwidth}{@{}p{0.9cm}YYY@{}}
\toprule
& \textbf{Duty} & \textbf{Observable trace} & \textbf{Falsifier}\\
\midrule
R1 & Bounded authority of knowledge and culture: state what is known, uncertain, or needs qualified
advice; culture or religion never closes a safety or consent question. & Uncertainty and role
boundaries stated in audited sessions; referral for out-of-scope questions. & A session in which a
cultural or religious claim is used to end a safety or consent question rather than to refer it.\\
R2 & A private channel for every person, at every stage, that the partner cannot see or control;
interpretation is provided by a trained, confidentiality-bound professional interpreter with
no prior relationship with either participant's family, never by a family member. & A routinely offered, safely designed private slot; uptake logged
without content. & A person who wanted a private slot did not receive one, or a family member
interpreted a safety-relevant exchange.\\
R3 & Revocable yes: pause, delay, refusal, or non-completion is a normal outcome, without penalty, at
any point on the line. & A pause, delay, or refusal recorded without an accompanying penalty or
exclusion. & A participant who paused, delayed, or refused was excluded, penalised, or pressured to
reverse the decision.\\
R4 & Stop--assess--refer--record: a trigger from any point on the line (self-check, private slot,
peer community, direct observation) stops joint work, moves to private assessment by a qualified,
competence-bounded assessor, refers to a qualified multidisciplinary service such as Thailand's
One-Stop Crisis Centre (OSCC) or an equivalent designated service, counts referral at connection, and
records the sequence; joint work does not resume where coercive control is confirmed. No couple
mediation continues once coercive control is confirmed. \emph{Proposed extension, not yet
accepted:} that assessor should additionally
have no stake in whether the wedding proceeds and no prior personal, family or congregational
relationship with either participant that a reasonable observer would see as compromising
independence, drawn where possible from outside the couple's own community network (analogy
evidence only; \cite{ref158,ref159,ref162,ref163}). If a qualified assessor cannot be reached
within a pre-specified maximum interval, or assessment is inconclusive, the default is continued
stop, not resumption; the case is escalated to a named safeguarding lead, not closed by facilitator
judgement. & Time from trigger to a qualified private assessment; referral
offered, connected, and recorded. & Joint work continued, or resumed, after a trigger without
a documented safe assessment; a case resumed or closed by facilitator judgement because an assessor
was unavailable, without safeguarding-lead sign-off.\\
R5 & Repairable records: every entry has a named owner, a review date, and is corrected when new
information arrives. & A review completed on or before its scheduled date, with a dated correction
where warranted. & A record that is never reviewed, or a correction that overwrites the original
entry rather than appending to it.\\
\bottomrule
\end{tabularx}
\end{table*}

If danger appears imminent, the routine R4 sequence is suspended in favour of the organisation's
emergency protocol: immediate safety planning and contact with emergency services or the OSCC
emergency line take precedence over scheduling confidential assessment. Facilitators are trained to
recognise the difference and are never the final decision-maker for either pathway. A named
safeguarding lead, independent of routine facilitation duty, owns supervision, escalation decisions
and incident review; the assessor judges the individual case, and the safeguarding lead owns the
system.

\subsubsection*{Layer 2 --- moderated peer community after the course (proposed extension)}
After delivery, a moderated peer space, anonymous by choice, is proposed as a route through which the
same rules keep operating rather than as a service in its own right: a place to keep the self-check
and the private channel alive, with a moderator trained to route a real trigger to R4 and never to
mediate. It carries no promise of support beyond that routing function. The nearest evidence families
are online peer-support research, where moderation quality is repeatedly identified in realist syntheses
as the factor that determines whether a space helps or harms \cite{ref140,ref144}, with mixed
effectiveness reported for youth mental-health peer forums specifically \cite{ref139}, and
technology-facilitated-abuse
research, which documents how an unmoderated or poorly bounded digital channel can itself become a
vector for surveillance, harassment, or coercive control \cite{ref149,ref150}. Neither literature was
generated on a Muslim or Thai premarital population, and the pairing is by analogy, not by direct
evidence for this layer. A pilot of Layer 2 requires, at minimum, technically enforced anonymity (no
real names, no shared-device login trace), a moderator-independent one-tap report channel that
routes directly to R4, and a periodic anonymous safety poll of members; none of this exists yet, and
Layer 2 should not be piloted without it.

\subsubsection*{One test over the whole line}
The programme theory has a single, pre-specified failure criterion, applied to the mechanism, the
rules, and every layer together rather than to any one point on it. Stated verbatim from the accepted
abstract: \tphe{} fails if it raises knowledge scores, attendance, satisfaction, or marriage completion
while leaving informed choice, recognition of coercion, reproductive and mental-health literacy, safe
help-seeking, and accountable referral unchanged. A polished, well-attended programme that leaves
those outcome domains flat has failed as public-health infrastructure, whatever its conventional
metrics show.

\subsection{The declared ethical frame}
The programme's own ethical language is named here as a declared frame, not as evidence: it states the
duty an actor is held to, while the evidence in Table~\ref{tab:rules} and elsewhere states what is
known about the corresponding construct. The two are never substituted for one another.
\emph{Faqr} (acknowledging one's own incompleteness and dependence) governs the humility required by
R1. \emph{Taqwa} (conscientious restraint in the use of knowledge and power) governs the same
boundary in R1 together with the stop obligation in R4. \emph{Amanah} (authority and knowledge held as
an entrusted responsibility, to be returned to whom it belongs) governs the private channel of R2, the
revocable consent of R3, and the repairable records of R5. \emph{Shura} (consultation when a limit is
reached) governs the assessor and referral step within R4. \emph{'Adl} (justice, tested by whether the
less powerful person keeps dignity, rights, voice, and a route to repair) is the standard the
failure criterion applies. Table~\ref{tab:ethic} pairs each term with the rule it governs, the
evidence category that speaks to the underlying construct, distinct from the term itself, and an
evidence tier (external evidence, analogy, open, or practitioner-derived) so a reader can see how
directly each pairing is supported. None of the humility- or compassion-adjacent literature cited below
was generated on a Muslim or Thai sample; it is adjacent-population evidence, not evidence for this
population.

\begin{table*}[t]
\centering\footnotesize
\caption{The declared ethical frame: term, gloss, governed rule, evidence category paired with it, and
the evidence tier. The term names a duty; the evidence category speaks to the construct, not to the
term. External-evidence and analogy citations are from adjacent, non-Thai, non-Muslim, non-premarital
populations unless stated otherwise. R4's no-stake-assessor clause, governed here by \emph{Shura},
is a proposed extension, not yet part of the accepted programme (Table~\ref{tab:rules}).}
\label{tab:ethic}
\begin{tabularx}{\textwidth}{@{}p{1.4cm}YYYp{1.7cm}@{}}
\toprule
\textbf{Term} & \textbf{Plain-language gloss} & \textbf{Governs} & \textbf{Evidence category paired with it} & \textbf{Evidence tier}\\
\midrule
Faqr & Acknowledging one's own incompleteness and dependence. & R1 & Intellectual and relational
humility, self-compassion \cite{ref209,ref211,ref212,ref213}; declared-frame companion literature in
Islamic medical ethics \cite{ref226,ref228}. & External evidence (analogy); no Muslim or Thai samples
for the humility\slash compassion constructs.\\
Taqwa & Conscientious restraint in the use of knowledge, speech, judgement, and power. & R1, R4 &
Self-regulation and provider-boundary and referral literature. & Open: no evidence pairing identified
for this construct at this population.\\
Amanah & Authority and knowledge held as an entrusted responsibility, returned to whom it belongs. &
R2, R3, R5 & Consent, autonomy, and accountability literature. & Open: no evidence pairing identified
for this construct at this population.\\
Shura & Consultation when a limit is reached. & R4 (assessor and referral step) & Independent-advocate
and dual-relationship literature \cite{ref158,ref159,ref162,ref163}. & Analogy: transplant-advocate and
dual-relationship ethics, not domain-matched.\\
'Adl & Justice: whether the less powerful person keeps dignity, rights, voice, and a route to
repair. & The failure criterion & Outcome-domain and safeguarding literature the failure criterion is
built from (see discussion \S4.3). & Practitioner-derived: the criterion itself is authored for this
paper, not drawn from a single cited instrument.\\
\bottomrule
\end{tabularx}
\end{table*}

\subsection{Labelling: what is accepted and what is proposed}
The mechanism, the five rules, the stop rule, and the failure criterion restate the accepted abstract
and carry its claim ceiling: no effectiveness claim is made anywhere in this line. \textbf{Layer 0
(signals from first contact and the anonymous self-check), Layer 2 (the moderated peer community), and
the no-stake assessor are extensions proposed for co-design after the abstract's acceptance.} They are
not part of the accepted programme theory, they add no outcome claim of their own, and any public text
using them must say so explicitly. Where the evidence base for a proposed extension is open in either
direction, that status is stated plainly rather than framed as a mild or roughly neutral uncertainty.

\subsection{Testable propositions}
Table~\ref{tab:props} converts the framework into a research agenda capable of producing disconfirming
evidence. These are propositions for prospective testing, not results.

\begin{table*}[t]
\centering\footnotesize
\caption{Testable propositions, competing explanations, and evidence that would weaken or falsify each.}
\label{tab:props}
\begin{tabularx}{\textwidth}{@{}p{0.6cm}YYY@{}}
\toprule
& \textbf{Proposition} & \textbf{Competing explanation} & \textbf{Weakening / falsifying evidence}\\
\midrule
P1 & A governance layer improves decision control beyond content-only education & Knowledge or
facilitator warmth explains any gain & No added gain in decision control after controlling for
knowledge and facilitator effects\\
P2 & A routine private channel makes materially relevant concerns more visible & Private contact adds
burden without information & No increase in safe voice or appropriate routing, with more distress,
retaliation, or privacy risk\\
P3 & R1's cultural-intelligibility component improves comprehension and trust only when culture or
religion cannot override consent or safety & Cultural tailoring alone suffices & Tailoring improves
satisfaction but leaves consent pressure or unsafe authority unchanged\\
P4 & A stop rule prevents routine joint education from worsening risk when fear or coercion appears &
Skilled facilitators can manage all concerns jointly & The stop rule causes more net harm than
governed alternatives, or does not reduce unsafe joint exposure\\
P5 & Referral quality depends on timeliness, accessibility, and follow-through, not offer alone &
Information is enough & Connection and safety do not vary by latency, accessibility, or
follow-through\\
P6 & Conventional programme success can coexist with \tphe{} failure & Conventional outcomes
represent benefit & Conventional outcomes reliably predict decision control, help-seeking, and
referral fidelity\\
P7 & R5 (repairable records) is necessary because relevant information emerges over time & One-time
assessment suffices & Review and correction add no safety or accuracy benefit and introduce more harm
or burden\\
\bottomrule
\end{tabularx}
\end{table*}
